\section{Nonvanishing at prime indices}\label{sec:nonvanishing}

This section proves that $T_n(r)\ne0$ for every rational $r$ and every large prime $n$
(\cref{thm:nonvanishing} with \cref{prop:local-dets}), the third input of the criterion, used in
\cref{sec:proofs}. \emph{Heuristic.} When $h$ and $K$ are fixed multiples of a prime $p$, every
residue class modulo $p$ contains the same number of poles. In a basis of polynomials adapted to the
classes, the distribution identity of \cref{sec:ordinary} makes the scaled matrix block triangular
modulo $p$, and each diagonal block reduces to one of three fixed rational matrices that do not
depend on $p$. If their determinants are nonzero, $\Delta(r)$ is a $p$-adic unit times an explicit
power of $p$. The numerator $a$ enters only through $Q=\sigma y^2=\sigma q^2/a^2$.

Let $(t,\gamma)\in\{(6,7),(4,5)\}$, and for a prime $p$ consider the determinant with
\[
  h=tp,\qquad K=\gamma p
\]
(the construction with $n=p$). Put $Q=\sigma y^2$, so that $w=1/Q$.

\subsection{The local matrices and the statement}\label{subsec:local-matrices}

For $s\in\{0,1\}$ let $\mathcal V_s$ be the $\Q$-vector space spanned by $\Q[z]$ and the
fractions $1/(z+l+\frac12)$, $s\le l<s+\gamma$, and define $\ell_Q\colon\mathcal V_s\to\Q$ by
\begin{equation}\label{eq:ellQ}
  \ell_Q(g)=\sum_{r\ge0}[z^r]g'\cdot M_r(1/Q)\quad(g\in\Q[z]),\qquad
  \ell_Q\Bigl(\frac1{z+l+\frac12}\Bigr)=4\sum_{i<l}\frac{Q^{\,l-i}}{(2i+1)^2},
\end{equation}
where $M_r(v)=\sum_{n\ge0}n^rv^n$ is the rational function of \cref{subsec:distribution}. For
$(s,e)\in\{(0,0),(1,0),(1,1)\}$ let $L_{s,e}(Q)$ be the $t\times t$ matrix with entries
\[
  \ell_Q\Bigl(\frac{z^{i+j}(z+1)^e}{\prod_{l=s}^{s+\gamma-1}(z+l+\frac12)}\Bigr),\qquad 0\le i,j<t .
\]
Its entries are rational numbers depending on $Q$ and $(t,\gamma)$ only.

\begin{theorem}\label{thm:nonvanishing}
Assume that the three determinants $\det L_{s,e}(Q)$ are nonzero. Then for every $r\in\Q$ and
every sufficiently large prime $p$ one has $\Delta(r)\ne0$, hence $T(r)\ne0$. More precisely,
$v_p(\Delta(r))=-pt(\gamma+2-t)+t$.
\end{theorem}

For $(t,\gamma)=(6,7)$ the valuation is $-18p+6$ and for $(4,5)$ it is $-12p+4$. We checked this
prediction against exact computations of $\Delta(0)$ at primes $p=13,17$ for five of the six
threshold rows, with $h=52,68,102$ (Appendix~\ref{app:consistency}).

\begin{proof}[Proof of \cref{thm:nonvanishing}]
Let $p$ be a prime not dividing $2aq(q^2-\sigma a^2)$ or the denominator of $r$, with $p>2\gamma+1$
(so $p^2>2K$), and not dividing the numerators and denominators of the three determinants or the
denominators of the entries of the $L_{s,e}(Q)$. All but finitely many primes qualify. We use the
notation of \cref{sec:ordinary}; every residue class $c\in\{1,\dots,p\}$ contains exactly
$\gamma$ poles.

\emph{Step 1: a basis adapted to the classes.} The ideals $((u-c)^t)$, $1\le c\le p$, of
$\Z_p[u]$ are pairwise comaximal, because $c-c'$ is a unit for $c\ne c'$. By the Chinese remainder
theorem there are $f_{c,i}\in\Z_p[u]$ of degree less than $tp$, for $1\le c\le p$ and $0\le
i<t$, with $f_{c,i}\equiv(u-c)^i\pmod{(u-c)^t}$ and $f_{c,i}\equiv0\pmod{(u-c')^t}$ for $c'\ne c$.
They form a $\Z_p$-basis of the polynomials of degree less than $h$ with $\Z_p$ coefficients, so
their coefficient matrix $C$ has $\det C\in\Z_p^\times$, and by \cref{lem:change-of-basis} the
Gram matrix $\mathbf G=[\mu_r(\varphi f_{c,i}f_{c',j})]$ has $\det\mathbf G=(\det C)^2\Delta(r)$.
Write $f_{c,i}(b+pz)=p^iF_{c,i,b}(z)$; then $F_{c,i,b}\in\Z_p[z]$, $F_{c,i,c}\equiv z^i$ and,
for $b\ne c$, $F_{c,i,b}\in p^{t-i}\Z_p[z]\subset p\Z_p[z]$.

\emph{Step 2: local weights.} By~\eqref{eq:local-weight}, for $b\ne p$,
$\varphi(b+pz)=p^{-\gamma}\Psi_b(z)$ with $\Psi_b=N_b/P_b$, $N_b\in\rest{\Z_p}$ and
\[
  N_b\equiv\lambda_b\pmod p,\qquad P_b(z)=\prod_{l\in J_b}\bigl(z+l+\tfrac12\bigr),
\]
where $\lambda_b=b\prod_{k\notin\Pi_b}(b-x_k)^{-1}$ is a unit. Here $J_b=\{0,\dots,\gamma-1\}$ if
$b\le\frac{p-1}2$ and $J_b=\{1,\dots,\gamma\}$ if $b\ge\frac{p+1}2$: indeed
$x_k-b=-(2k+2b+1)/2$, and $2k+2b+1=(2l+1)p$ for $0\le k<\gamma p$ gives exactly these ranges of
$l$, with $z_k=-l-\frac12$ \emph{exactly}. For $b=p$ the numerator $u=p(1+z)$ supplies one more
factor $p$: $\varphi(p+pz)=p^{1-\gamma}\Psi_p(z)$ with $N_p\equiv\lambda_p(1+z)$, where now $\lambda_p=\prod_{k\notin\Pi_p}(-x_k)^{-1}$ is a unit.

\emph{Step 3: the scaled matrix.} By \cref{lem:distribution},
$\mathbf G_{(c,i),(c',j)}=\sum_bw^bp^{\,i+j-1-\gamma}\ell_b(\tilde\varphi_bF_{c,i,b}F_{c',j,b})$, where
$\tilde\varphi_b=\Psi_b$ for $b\ne p$ and $\tilde\varphi_p=p\Psi_p$. Let $\mathbf M$ be obtained from
$\mathbf G$ by multiplying row $(c,i)$ by $p^{\gamma+1-i}$, and moreover by $p^{-1}$ if $c=p$, and
column $(c',j)$ by $p^{-j}$. Then
\[
  \mathbf M_{(c,i),(c',j)}=p^{-\one_{c=p}}\sum_{b=1}^pw^b\,\ell_b\bigl(\tilde\varphi_bF_{c,i,b}F_{c',j,b}\bigr),
\]
which is $p$-integral by \cref{lem:local-integrality}(b): for $c\ne p$ directly; for $c=p$ because
$F_{p,i,b}\in p\Z_p[z]$ when $b\ne p$, and $\tilde\varphi_p\in p\,V_p$.

\emph{Step 4: block triangularity.} Modulo $p$, $F_{c,i,b}\equiv\one_{c=b}z^i$, and $\ell_b$ is
$\Z_p$-linear on the lattice of \cref{lem:local-integrality}(b), so a term with $b\ne c$ or $b\ne c'$
vanishes modulo $p$ unless it is divided by $p$. Hence for $c\ne p$ the row $(c,i)$ of
$\mathbf M$ is, modulo $p$, supported in the columns $(c,j)$, with entries
$w^c\ell_c(\Psi_cz^{i+j})$. For $c=p$, the entries in columns $(p,j)$ receive from $b\ne p$ terms
in $p^{2t-i-j-1}\Z_p\subset p\Z_p$, so they are $w^p\ell_p(\Psi_pz^{i+j})$ modulo $p$. Ordering the
classes with $p$ last, $\mathbf M$ is block lower triangular modulo $p$, and
\begin{equation}\label{eq:det-product}
  \det\mathbf M\equiv\prod_{c=1}^{p}\det\bigl[w^c\ell_c(\Psi_cz^{i+j})\bigr]_{i,j<t}\pmod p .
\end{equation}

\emph{Step 5: reduction of the local functionals.} We claim that for $N\in\rest{\Z_p}$,
$\ell_c(N/P_c)\equiv\ell_Q(\bar N/P_c)\pmod p$, where $\bar N$ is the polynomial obtained by
reducing the coefficients of $N$ and $\ell_Q$ is reduced modulo $p$. For the restricted-series
part this follows from $v=w^p\equiv w=1/Q\pmod p$ and the fact that $M_r$ has denominator
$(1-v)^{r+1}$, a unit. For a pole $k\in\Pi_c$ with $z_k=-l-\frac12$, the sum over $b\ne c$
in~\eqref{eq:ell-pole} is divisible by $p^2$, and $p^2X\equiv0$, so
$\ell_c(1/(z-z_k))\equiv4\sigma^ky^{2k+1}w^{-c}p^2\tau_k$. In $p^2\tau_k$ only the terms with
$p\mid2i'+1$ survive; writing $2i'+1=(2i+1)p$, they are those with $i<l$ (because
$2k+1=(2l+1)p-2c$ and $1\le c\le p$), and they contribute
$\sigma^{ip+(p-1)/2}y^{-(2i+1)p}/(2i+1)^2$. Using $\sigma^p=\sigma$, $y^p\equiv y$, $w^{-c}=\sigma^c
y^{2c}$, $y^{2k+2c+1}=y^{(2l+1)p}$ and $k+c+\frac{p-1}2=(l+1)p-1\equiv l\pmod2$, we obtain
\[
  \ell_c\Bigl(\frac{1}{z-z_k}\Bigr)\equiv4\sum_{i<l}\sigma^{l+i}y^{2l-2i}(2i+1)^{-2}
  =4\sum_{i<l}\frac{Q^{\,l-i}}{(2i+1)^2}\pmod p,
\]
which is~\eqref{eq:ellQ}. Here $y^p\equiv y$ means congruence in $\Z_p$ modulo $p$ after
reduction, $y=q/a$ being a $p$-adic unit, and the reduction of rational numbers to $\mathbb F_p$
is the one used throughout. Since the decomposition of $N/P_c$ into a restricted series and
principal parts has integral coefficients (\cref{lem:local-integrality}(b)) and commutes with
reduction, the claim follows.

\emph{Step 6: conclusion.} By Steps 2 and 5 the $c$-th block in~\eqref{eq:det-product} is
congruent to $w^c\lambda_c$ times $L_{s_c,e_c}(Q)$, with $s_c=\one_{c>(p-1)/2}$ and $e_c=\one_{c=p}$.
Hence $\det\mathbf M$ is a $p$-adic unit. As $\det\mathbf M=p^{\,pt(\gamma+1)-pt(t-1)-t}\det\mathbf
G$ and $\det C$ is a unit, $v_p(\Delta(r))=-pt(\gamma+2-t)+t$.
\end{proof}

Krisman's Lemma~P \cite[Thm.~8.1]{Krisman2026chi2} follows the same pattern with $h=3\ell$,
$K=h$, the Lagrange basis at the poles and three $3\times3$ blocks. With $K>h$ there is no
Lagrange basis at the poles, which is why we use the Taylor jets of Step~1.

\subsection{The local determinants}\label{subsec:local-dets}

\begin{proposition}\label{prop:local-dets}
For $(t,\gamma)=(6,7)$ and $(s,e)=(0,0),(1,0),(1,1)$,
\[
  (Q-1)^{24}\det L_{s,e}(Q)=c_{s,e}\,Q^{\,a_{s,e}}(Q-1)^{b_{s,e}}\,\Pi_{s,e}(Q)
\]
with nonzero rationals $c_{s,e}$, $(a_{s,e},b_{s,e})=(9,20),(11,20),(13,18)$, and primitive
integer polynomials $\Pi_{s,e}$ of degrees $16,20,20$. For $(t,\gamma)=(4,5)$ the same holds with
$(Q-1)^{16}$, $(a_{s,e},b_{s,e})=(5,15),(6,15),(7,14)$ and degrees $6,9,9$. In all six cases the
roots of $\Pi_{s,e}$ have modulus less than Cauchy's bound $1+\max_i\abs{\pi_i/\pi_{\deg}}$, which is less than
$23.19$, $136.73$, $188.09$ for $(6,7)$ and less than $1.71$, $9.38$, $5.08$ for $(4,5)$. In
particular $\det L_{s,e}(Q)\ne0$ whenever $\abs Q\ge189$, hence whenever $q\ge14a$.
\end{proposition}

\begin{proof}
The entries of $(Q-1)^4L_{s,e}(Q)$ are polynomials in $Q$: the derivative of the polynomial part has
degree at most $3$, so only $M_0,\dots,M_3$ occur, and the pole values are polynomials. The
determinants were computed exactly, twice and independently (Appendix~\ref{app:local-certificates}): once
symbolically in $Q$, and once by exact interpolation at rational values, with agreement. The bound on
the roots is Cauchy's, evaluated exactly. Finally $q\ge14a$ gives $\abs Q=q^2/a^2\ge196>189$.
\end{proof}

The proposition is used with $\abs Q\ge196$; all six threshold rows have $q/a\ge43$.
